\documentclass[11pt,a4paper]{article}
\usepackage[margin=19mm]{geometry}
\usepackage{amsmath,amssymb,amsthm,hyperref}
\newtheorem{theorem}{Theorem}
\newtheorem{lemma}[theorem]{Lemma}
\newcommand{\gen}[1]{\langle #1\rangle}
\newcommand{\Aut}{\operatorname{Aut}}
\title{Stabilizers in geometric halving}
\author{Human Compression project}
\date{}
\begin{document}
\maketitle

We give a stabilizer proof of the coset identification in
Theorems~B and~4.4 of~\cite{PT}. The proof does not
require the halved geometry to be residually connected. This distinction
matters: the assertion about residual connectedness in the inspected preprint
has counterexamples, whereas the coset identification below remains
valid. No claim is made about an uninspected later edition.

\section{The construction and the retained geometric facts}
Let $\Gamma$ be a regular hypertope of finite rank $n\ge3$ with
distinguished generators
\[
 G=\Aut(\Gamma)=\gen{a,b,r_2,\ldots,r_{n-1}},
 \qquad [a,r_i]=1\quad(i\ge2).
\]
The generators are involutions and change, respectively, the objects
of types $0,1,2,\ldots,n-1$ in the base chamber
$(v,\{v,w\},x_2,\ldots,x_{n-1})$, where $w=av$.
Assume that the vertex--edge truncation is a simple graph $Q$, and
that an edge is incident with a higher object exactly when both its
endpoints are incident with that object. These are conditions
(B1) and (B2) of~\cite{PT}.

For a graph $S$, its bipartite double cover $D(S)$ has vertices
$(u,\epsilon)$, with $\epsilon\in\{0,1\}$, and edges
$(u,0)(z,1)$ whenever $uz$ is an edge of $S$.
Choose the component $C$ of $D(Q)$ containing $(v,0)$.
For a higher object $x$, let $Q_x$ be its vertex--edge shadow.
The objects of the halved incidence system $L$ are the vertices
of $C$, with their two types, and the pairs $(x,C_x)$, where $C_x$
is a component of $D(Q_x)$ contained in $C$.
Vertex incidence is graph adjacency or membership in $C_x$;
two higher objects are incident when their original objects are
incident and their chosen components intersect.

We retain the previously proved component-intersection lemma for
this construction. It uses connectedness of the original flag
shadows, (B1), and (B2), and does not assume residual connectedness
of $L$. Its consequences needed here are that $L$ is a thin geometry
and that its chambers correspond bijectively to the original
chambers whose oriented edge has its lift in $C$. The component
of each higher object is uniquely determined by the lifted first
endpoint. In particular, when $Q$ is nonbipartite every original
chamber occurs; when $Q$ is bipartite exactly the chambers whose
first endpoint belongs to the chosen color class occur.
The full retained proof, including extension of every flag to a
chamber, is charged as a dependency of this note.

\section{An elementary parity lemma}
\begin{lemma}\label{parity}
Suppose a group $K=\gen{a,b,R}$ acts transitively on the vertices
and edges of a connected graph $S$, the stabilizer of a base vertex
$u$ is $P=\gen{b,R}$, and the edge orbit is that of $\{u,au\}$.
Suppose $a,b$ are involutions and $a$ commutes with every member
of $R$. Put $K^+=\gen{aba,b,R}$.
Then $K^+$ is normal in $K$ with index at most two.
It is proper exactly when $S$ is bipartite. In that case it is
the subgroup preserving the two color classes. Equivalently,
it is the stabilizer of either component of $D(S)$ under the
action preserving the second coordinate.
\end{lemma}
\begin{proof}
Conjugation by $a$ exchanges $b$ and $aba$ and fixes $R$.
Thus $a$ normalizes $K^+$, and $K=K^+\gen a$. If $K^+$ is proper,
the character $K\longrightarrow K/K^+\cong C_2$ is constant on
cosets of $P$, since $P\le K^+$. It colors the vertices $gu$
by the coset of $g$. The endpoints $gu$ and $gau$ have different
colors, so the coloring is a bipartition of $S$.

Conversely, the bipartition of a connected bipartite graph is
unique up to exchange. The action of $K$ on its two classes is
a homomorphism to $C_2$. Every element of $P$ fixes $u$ and
therefore preserves the classes, while $a$ exchanges them.
Consequently its kernel contains $K^+$ and excludes $a$.
The index bound gives equality with the kernel.
Finally, an even walk stays in a color class and an odd walk
changes it. Hence $D(S)$ has two components in the bipartite
case, and one otherwise; their stabilizers have exactly the
description asserted.
\end{proof}

\section{All maximal parabolic stabilizers}
Put
\[
 t_0=aba,\quad t_1=b,\quad t_i=r_i\ (i\ge2),
 \qquad H=\gen{t_0,\ldots,t_{n-1}},\qquad
 H_i=\gen{t_j:j\ne i}.
\]
\begin{theorem}\label{cosets}
The group $H$ acts regularly on the chambers of $L$ and is its
full group of type-preserving automorphisms. In its base chamber,
the stabilizer of the object of type $i$ is exactly $H_i$.
Consequently $L$ is isomorphic to the typed coset incidence
geometry of $H$ and $(H_i)_{0\le i<n}$.
None of these assertions requires residual connectedness of $L$.
\end{theorem}
\begin{proof}
Apply Lemma~\ref{parity} first to $G$ and $Q$. The original
vertex stabilizer is $P=\gen{b,r_2,\ldots,r_{n-1}}$; this is
the standard parabolic stabilizer in a regular hypertope.
If $Q$ is nonbipartite, then $H=G$. If it is bipartite, $H$
has index two and preserves the chosen class. The chamber
correspondence above and the regular action of $G$ on original
chambers show that $H$ acts regularly on the chambers of $L$.

The stabilizer of $(v,0)$ is $P=H_0$. The stabilizer of $(w,1)$
is $aPa=\gen{aba,r_2,\ldots,r_{n-1}}=H_1$.
Both subgroups are contained in $H$, so these are also the
stabilizers for the restricted action.

Fix $i\ge2$ and set
\[
 K_i=\gen{a,b,r_j: j\ge2,\ j\ne i}.
\]
This is the full stabilizer of $x_i$ in $G$. In the original
residue at $x_i$, it is transitive on vertices and edges, and
its vertex stabilizer is $\gen{b,r_j:j\ge2,\ j\ne i}$.
The shadow graph $Q_{x_i}$ is connected: every truncation on
at least two types of a residually connected geometry is
connected. Its edges are the translates of the base edge.
Thus Lemma~\ref{parity} applies inside $K_i$, even if its action
on the shadow graph has a kernel. It says that the stabilizer
in $K_i$ of the chosen component of $D(Q_{x_i})$ is precisely
\[
 K_i^+=\gen{aba,b,r_j:j\ge2,\ j\ne i}=H_i.
\]
An element fixing $(x_i,C_i)$ must fix $x_i$ and preserve $C_i$.
Since $H_i\le H$, its stabilizer in $H$ is exactly $H_i$.

The elements $t_i$ carry the base chamber to its unique
$i$-adjacent chamber. For $i\ge2$ this follows from the old
adjacency and the uniqueness of the chosen components through
the retained endpoint. The element $b$ fixes $v$ and changes
the other endpoint of the base edge, so it gives adjacency
of type $1$. Conjugating this operation by $a$ gives $aba$,
which fixes $w$ and changes the endpoint of type $0$.
It preserves every higher component, because it fixes the
corresponding higher object and the retained lifted endpoint.
These are genuine changes, by thinness of the original panels.
Since the $t_i$ generate $H$ and $H$ is chamber-transitive,
the chamber graph of $L$ is connected. In a thin geometry,
a type-preserving automorphism fixing a chamber fixes each
adjacent chamber; induction along galleries shows that it
fixes every object. Thus the action of the full automorphism
group on chambers is free. The regular subgroup $H$ is already
transitive, and hence is the full automorphism group.

Lastly send a left coset $hH_i$ to the image under $h$ of the
base object of type $i$. The stabilizer calculation makes
this a bijection on each type. Intersecting cosets give
incident objects. Conversely, incident objects extend to a
chamber of $L$, and chamber transitivity supplies an element
in the two corresponding cosets. This proves the incidence
isomorphism, with distinct types retained even if subgroups
coincide.
\end{proof}

The theorem proves the literal coset assertion of~\cite{PT};
it does not infer the subgroup intersection property from
chamber regularity. A stabilizer of a flag of two or more
higher objects need not be generated by the remaining $t_i$.
That is precisely where the previously established parity
condition and the counterexamples to residual connectedness
enter. The present argument supplies a missing proof of a
known assertion, not a newly solved open conjecture.

\begin{thebibliography}{1}
\bibitem{PT} C.~A. Piedade and P. Tranchida,
\emph{Constructing new geometries: a generalized approach to
halving for hypertopes}, arXiv:2405.19050v1 (29 May 2024),
Theorems B and 4.4, Proposition 4.3.
\end{thebibliography}
\end{document}
